% =====================================================================
\chapter{Estimación de parámetros}
\label{tema:4}
% =====================================================================

Entrenar un modelo es elegir sus parámetros a partir de datos. Este tema presenta los dos enfoques clásicos de la estimación: el frecuentista, cuyo método principal es la máxima verosimilitud, y el bayesiano, que trata los parámetros como variables aleatorias con una distribución a priori. Se sigue la §9.1 y las §8.1–8.2 de~\cite{bt}, las lecciones~21 a~23 de~\cite{mit6041} y las clases~10 a~15 de~\cite{mit1805}; la conexión con el aprendizaje automático procede de~\cite[§6.6 y §8.3]{mml} y de~\cite[§2.1–2.3]{bishop}.

\section{Estimadores, sesgo y error cuadrático medio}

\begin{definicion}[Modelo estadístico y estimador]
Un \emph{modelo estadístico} es una familia de distribuciones $\{p(\cdot\,;\theta):\theta\in\Theta\}$ indexada por un parámetro $\theta$. Dada una muestra $X_1,\dots,X_n$ \iid\ con distribución $p(\cdot\,;\theta)$ para un $\theta$ desconocido, un \emph{estimador} es una función $\hat\theta_n=g(X_1,\dots,X_n)$ de la muestra. Es una variable aleatoria.
\end{definicion}

\begin{definicion}[Sesgo y error cuadrático medio]
El \emph{sesgo} de $\hat\theta_n$ es $b(\hat\theta_n)=\E[\hat\theta_n]-\theta$; el estimador es \emph{insesgado} si $b=0$ para todo $\theta$. Su \emph{error cuadrático medio} es $\mathrm{ECM}(\hat\theta_n)=\E\bigl[(\hat\theta_n-\theta)^{2}\bigr]$~\cite[§9.1]{bt}.
\end{definicion}

\begin{proposicion}[Descomposición sesgo--varianza]\label{prop:ecm}
$\mathrm{ECM}(\hat\theta_n)=b(\hat\theta_n)^{2}+\Var(\hat\theta_n)$.
\end{proposicion}

\begin{proof}
Sea $m=\E[\hat\theta_n]$. Entonces $\hat\theta_n-\theta=(\hat\theta_n-m)+(m-\theta)$. Al elevar al cuadrado y tomar esperanzas, el término cruzado es $2(m-\theta)\E[\hat\theta_n-m]=0$, y quedan $\Var(\hat\theta_n)$ y $(m-\theta)^{2}=b^{2}$.
\end{proof}

Por el \cref{cor:media}, la media muestral $\Xbar_n$ es un estimador insesgado de la media, con $\mathrm{ECM}=\sigma^{2}/n$.

\section{Máxima verosimilitud}

\begin{definicion}[Estimador de máxima verosimilitud]\label{def:emv}
Dada una realización $x_1,\dots,x_n$ de una muestra \iid, la \emph{verosimilitud} es $\mathcal{L}(\theta)=\prod_{i=1}^{n}p(x_i;\theta)$ y la \emph{log-verosimilitud} es $\ell(\theta)=\sum_{i=1}^{n}\ln p(x_i;\theta)$. El \emph{estimador de máxima verosimilitud} (EMV) es
\[
\hat\theta_{\mathrm{EMV}}=\argmax_{\theta\in\Theta}\ell(\theta)=\argmin_{\theta\in\Theta}\Bigl(-\frac1n\sum_{i=1}^{n}\ln p(x_i;\theta)\Bigr).
\]
\end{definicion}

Las dos formas son equivalentes porque $\ln$ es creciente y dividir por $-n$ invierte el sentido de la optimización. La segunda es la log-verosimilitud negativa media, que el módulo~1 identificó con la entropía cruzada entre la distribución empírica de los datos y el modelo. En consecuencia, \emph{preentrenar un LLM minimizando la entropía cruzada es calcular un estimador de máxima verosimilitud}, con $\theta$ de miles de millones de componentes~\cite[§8.3]{mml}.

\begin{ejemplo}[Bernoulli]\label{ej:emv-bernoulli}
Con $k$ éxitos en $n$ ensayos de $\Ber(p)$, $\ell(p)=k\ln p+(n-k)\ln(1-p)$. Derivando, $\ell'(p)=k/p-(n-k)/(1-p)=0$ da $\hat p=k/n$, y $\ell''(p)=-k/p^{2}-(n-k)/(1-p)^{2}<0$ confirma el máximo. Con 7 caras en 10 lanzamientos, $\hat p=0{,}7$ y $\ell(0{,}7)=-6{,}109$, mayor que $\ell(0{,}5)=10\ln 0{,}5=-6{,}931$.
\end{ejemplo}

Para la categórica con conteos $N_1,\dots,N_K$, el EMV es la frecuencia relativa $\hat p_j=N_j/\sum_k N_k$; es el teorema del EMV del modelo de bigramas del módulo~1, demostrado allí con multiplicadores de Lagrange. Contar y normalizar es, por tanto, máxima verosimilitud.

\section{El EMV de la normal}

\begin{teorema}[EMV de la normal]\label{teo:emv-normal}
Sea $x_1,\dots,x_n$ una muestra de $\N(\mu,\sigma^{2})$ con $n\ge 2$ y no todos los $x_i$ iguales. El EMV de $(\mu,\sigma^{2})$ es
\[
\hat\mu=\bar x=\frac1n\sum_{i=1}^{n}x_i,\qquad \hat\sigma^{2}=\frac1n\sum_{i=1}^{n}(x_i-\bar x)^{2}.
\]
\end{teorema}

\begin{proof}
Con $v=\sigma^{2}$, la log-verosimilitud es
\[
\ell(\mu,v)=-\frac n2\ln(2\pi v)-\frac{1}{2v}\sum_{i=1}^{n}(x_i-\mu)^{2}.
\]
\emph{Paso 1.} $\partial\ell/\partial\mu=\frac1v\sum_i(x_i-\mu)=0$ da $\mu=\bar x$, para cualquier $v>0$.

\emph{Paso 2.} Sea $S=\sum_i(x_i-\bar x)^{2}>0$. Entonces $\partial\ell/\partial v=-\frac{n}{2v}+\frac{S}{2v^{2}}=0$ da $v=S/n$.

\emph{Paso 3.} En $(\bar x,S/n)$ las segundas derivadas son $\partial^{2}\ell/\partial\mu^{2}=-n/v$, $\partial^{2}\ell/\partial\mu\,\partial v=-\frac{1}{v^{2}}\sum_i(x_i-\bar x)=0$ y $\partial^{2}\ell/\partial v^{2}=\frac{n}{2v^{2}}-\frac{S}{v^{3}}=-\frac{n}{2v^{2}}$. La hessiana es diagonal con entradas negativas, luego el punto es un máximo local; como es el único punto crítico y $\ell\to-\infty$ cuando $v\to 0^{+}$, $v\to\infty$ o $\abs\mu\to\infty$, es el máximo global.
\end{proof}

\begin{ejemplo}\label{ej:emv-normal}
Para los datos $(2,4,4,4,5,5,7,9)$, $\hat\mu=40/8=5$. Las desviaciones al cuadrado son $9,1,1,1,0,0,4,16$, que suman~32, y $\hat\sigma^{2}=32/8=4$, $\hat\sigma=2$.
\end{ejemplo}

\begin{proposicion}[Sesgo de la varianza muestral]\label{prop:sesgo}
Si $X_1,\dots,X_n$ son \iid\ con varianza $\sigma^{2}$, entonces $\E[\hat\sigma^{2}]=\frac{n-1}{n}\sigma^{2}$. En consecuencia, $s^{2}=\frac{1}{n-1}\sum_i(X_i-\Xbar_n)^{2}$ es un estimador insesgado de $\sigma^{2}$.
\end{proposicion}

\begin{proof}
Se usa la identidad $\sum_i(X_i-\Xbar_n)^{2}=\sum_i(X_i-\mu)^{2}-n(\Xbar_n-\mu)^{2}$, cuya deducción detallada se pide en el \cref{eje:4-sesgo}. La esperanza del primer término es $n\sigma^{2}$; la del segundo, $n\Var(\Xbar_n)=\sigma^{2}$ por el \cref{cor:media}. Por tanto, $\E\bigl[\sum_i(X_i-\Xbar_n)^{2}\bigr]=(n-1)\sigma^{2}$.
\end{proof}

En el \cref{ej:emv-normal}, la versión insesgada es $s^{2}=32/7\approx 4{,}571$. El sesgo se debe a que $\bar x$ se ajusta a los propios datos y las desviaciones respecto de él son, en media, menores que respecto de $\mu$. El EMV no es siempre insesgado, pero su sesgo tiende a~0 cuando $n\to\infty$.

\begin{proposicion}[Mínimos cuadrados como máxima verosimilitud]\label{prop:mco}
Supóngase $y_i=f_\theta(x_i)+\varepsilon_i$ con $\varepsilon_i\sim\N(0,\sigma^{2})$ independientes y $\sigma$ fija. Maximizar la verosimilitud de $(y_1,\dots,y_n)$ en $\theta$ equivale a minimizar el error cuadrático $\sum_i\bigl(y_i-f_\theta(x_i)\bigr)^{2}$.
\end{proposicion}

\begin{proof}
Cada $y_i\sim\N(f_\theta(x_i),\sigma^{2})$, de modo que $-\ell(\theta)=\frac n2\ln(2\pi\sigma^{2})+\frac{1}{2\sigma^{2}}\sum_i\bigl(y_i-f_\theta(x_i)\bigr)^{2}$. El primer sumando y el factor $1/(2\sigma^{2})>0$ no dependen de $\theta$.
\end{proof}

El módulo~8 usa la \cref{prop:mco} para justificar la regresión lineal; la entropía cruzada de la clasificación es su análogo para la categórica.

\section{Inferencia bayesiana y distribuciones conjugadas}

El EMV de la categórica asigna probabilidad nula a toda categoría no observada, y la pérdida de un texto que la contenga es infinita (módulo~1). El enfoque bayesiano lo remedia tratando el parámetro como una variable aleatoria con una distribución \emph{a priori} $\pi(\theta)$ y actualizándola con la regla de Bayes~\cite[§8.1]{bt}:
\[
\pi(\theta\mid x_{1:n})=\frac{\mathcal{L}(\theta)\,\pi(\theta)}{\int\mathcal{L}(\theta')\,\pi(\theta')\,d\theta'}\ \propto\ \mathcal{L}(\theta)\,\pi(\theta).
\]

\begin{definicion}[Estimadores bayesianos]
La \emph{media a posteriori} es $\E[\theta\mid x_{1:n}]$. El estimador \emph{máximo a posteriori} (MAP) es $\hat\theta_{\mathrm{MAP}}=\argmax_\theta\pi(\theta\mid x_{1:n})=\argmax_\theta\bigl(\ell(\theta)+\ln\pi(\theta)\bigr)$~\cite[§8.2]{bt}.
\end{definicion}

\begin{definicion}[Distribución beta]
$p\sim\BetaD(a,b)$, con $a,b>0$, si su densidad en $(0,1)$ es proporcional a $p^{a-1}(1-p)^{b-1}$. Su media es $a/(a+b)$ y, si $a,b>1$, su moda es $(a-1)/(a+b-2)$~\cite[§2.1.1]{bishop}.
\end{definicion}

\begin{teorema}[Conjugación beta--Bernoulli]\label{teo:beta}
Si $p\sim\BetaD(a,b)$ a priori y se observan $k$ éxitos en $n$ ensayos de Bernoulli independientes dado $p$, la distribución a posteriori es $\BetaD(a+k,\,b+n-k)$. En particular,
\[
\E[p\mid\text{datos}]=\frac{k+a}{n+a+b}.
\]
\end{teorema}

\begin{proof}
La verosimilitud es proporcional a $p^{k}(1-p)^{n-k}$. Multiplicando por la densidad a priori, la densidad a posteriori es proporcional a $p^{k+a-1}(1-p)^{n-k+b-1}$, que es el núcleo de una $\BetaD(a+k,b+n-k)$; la constante queda fijada por la normalización. La media se obtiene de la fórmula $a/(a+b)$ con los nuevos parámetros.
\end{proof}

Se dice que la beta es \emph{conjugada} de la Bernoulli: la distribución a posteriori pertenece a la misma familia que la a priori, y los datos se limitan a sumar conteos a los parámetros, que actúan como \emph{pseudoconteos}~\cite[clase~15]{mit1805}. El mismo argumento vale para $K$ categorías con la distribución de Dirichlet, de densidad proporcional a $\prod_{j=1}^{K}p_j^{\alpha_j-1}$ sobre el símplex: con conteos $N_1,\dots,N_K$, la distribución a posteriori es $\Dir(\alpha_1+N_1,\dots,\alpha_K+N_K)$~\cite[§2.2.1]{bishop}, \cite[§6.6]{mml}.

\begin{corolario}[Suavizado de Laplace como estimación bayesiana]\label{cor:laplace}
Con distribución a priori $\Dir(\alpha,\dots,\alpha)$ y conteos $N_1,\dots,N_K$ de total $N$,
\[
\E[p_j\mid\text{datos}]=\frac{N_j+\alpha}{N+K\alpha},\qquad
\hat p_{j,\mathrm{MAP}}=\frac{N_j+\alpha-1}{N+K(\alpha-1)}\quad(\alpha>1).
\]
El suavizado de Laplace del módulo~1 es la media a posteriori con $\alpha=1$ (a priori uniforme sobre el símplex) y también el MAP con $\alpha=2$.
\end{corolario}

\begin{table}[h]
\centering\small
\caption{Estimadores de una categórica con conteos $N=(3,1,0)$, $K=3$.}\label{tab:estimadores}
\begin{tabular}{@{}lll@{}}
\toprule
Estimador & Fórmula para $p_j$ & Resultado \\
\midrule
EMV & $N_j/N$ & $(0{,}75;\ 0{,}25;\ 0)$ \\
Media a posteriori, $\alpha=1$ & $(N_j+1)/(N+K)$ & $(4/7;\ 2/7;\ 1/7)\approx(0{,}571;\ 0{,}286;\ 0{,}143)$ \\
MAP, $\alpha=2$ & $(N_j+1)/(N+K)$ & igual que la fila anterior \\
\bottomrule
\end{tabular}
\end{table}

\begin{observacion}[Conexión con los LLM]
En una red neuronal no se suele fijar una distribución a priori explícita, pero el término $\ln\pi(\theta)$ del MAP reaparece como \emph{regularización}: con una a priori normal $\theta\sim\N(0,\tau^{2}I)$, $\ln\pi(\theta)=-\norm{\theta}^{2}/(2\tau^{2})+\text{cte}$, y maximizar $\ell(\theta)+\ln\pi(\theta)$ es minimizar la entropía cruzada más una penalización cuadrática de los pesos (\emph{weight decay}), que se estudia en el módulo~7~\cite[§8.3]{mml}.
\end{observacion}

\begin{lecturas}
\item Bertsekas y Tsitsiklis~\cite{bt}, §8.1 «Bayesian Inference and the Posterior Distribution», §8.2 «Point Estimation, Hypothesis Testing, and the MAP Rule» y §9.1 «Classical Parameter Estimation».
\item MIT 6.041SC~\cite{mit6041}, unidad~IV, lecciones~21–22 «Bayesian Statistical Inference» y~23 «Classical Statistical Inference~-~I».
\item MIT 18.05~\cite{mit1805}, clases~10b «Maximum Likelihood Estimates», 11 «Bayesian Updating with Discrete Priors», 12a «Bayesian Updating: Probabilistic Prediction» y 15 «Conjugate priors: Beta and normal».
\item Deisenroth, Faisal y Ong~\cite{mml}, §6.6 «Conjugacy and the Exponential Family» y §8.3 «Parameter Estimation».
\item Bishop~\cite{bishop}, §2.1.1 «The beta distribution», §2.2.1 «The Dirichlet distribution» y §2.3.4 «Maximum likelihood for the Gaussian».
\end{lecturas}

\section*{Ejercicios del tema 4}
\addcontentsline{toc}{section}{Ejercicios del tema 4}

\begin{ejercicio}{1}{eje:4-ceros}
Explique por qué el EMV de una categórica es problemático cuando hay pocos datos y cómo lo resuelve una distribución a priori de Dirichlet. Con conteos $(3,1,0)$, ¿qué probabilidad asigna al tercer valor el suavizado de Laplace?
\end{ejercicio}

\begin{ejercicio}{2}{eje:4-datos}
Para los datos $(1,3,5,7)$, supuestos procedentes de una normal, calcule el EMV de $\mu$ y de $\sigma^{2}$, y la varianza muestral insesgada $s^{2}$.
\end{ejercicio}

\begin{ejercicio}{2}{eje:4-moneda}
En 10 lanzamientos de una moneda salen 7 caras. Calcule $\hat p_{\mathrm{EMV}}$, $\ell(0{,}7)$ y $\ell(0{,}5)$, y la estimación de Laplace $(k+1)/(n+2)$. Interprete la diferencia entre ambas estimaciones.
\end{ejercicio}

\begin{ejercicio}{2}{eje:4-u2}
(a) Una caja contiene tres monedas con probabilidad de cara $1/3$, $1/2$ y $2/3$. Se extrae una y, en 80 lanzamientos, salen 49 caras. ¿Qué moneda maximiza la verosimilitud? (b) Si $p$ puede tomar cualquier valor de $[0,1]$, ¿cuál es el EMV? (c) Las duraciones de cinco bombillas, supuestas exponenciales de parámetro $\lambda$, son 2, 3, 1, 3 y 4 años. Obtenga el EMV de $\lambda$.
\fuente{adaptado de MIT 18.05~\cite{mit1805}, clase~10, \emph{Board Questions}~1 y~2.}
\end{ejercicio}

\begin{ejercicio}{2}{eje:4-u3}
Un cajón contiene tres monedas con probabilidad de cara $0{,}5$, $0{,}6$ y $0{,}9$. Se elige una al azar, se lanza 5 veces y sale una sola cara. Calcule la probabilidad a posteriori de cada moneda y la probabilidad predictiva de cara en el lanzamiento siguiente, antes y después de observar los datos.
\fuente{adaptado de MIT 18.05~\cite{mit1805}, clase~12, problema~1.}
\end{ejercicio}

\begin{ejercicio}{3}{eje:4-sesgo}
Sean $X_1,\dots,X_n$ independientes con media $\mu$ y varianza $\sigma^{2}$. Demuestre la identidad $\sum_i(X_i-\Xbar_n)^{2}=\sum_i(X_i-\mu)^{2}-n(\Xbar_n-\mu)^{2}$ y complete la demostración de la \cref{prop:sesgo}. Compruebe el caso $n=1$.
\end{ejercicio}

\begin{ejercicio}{3}{eje:4-beta}
Para $K=2$, la distribución a priori de Dirichlet simétrica es $\BetaD(\alpha,\alpha)$. Demuestre, sin usar el \cref{teo:beta}, que tras $k$ éxitos en $n$ ensayos la media a posteriori es $(k+\alpha)/(n+2\alpha)$. Puede usar que la media de una $\BetaD(a,b)$ es $a/(a+b)$. Compruebe los casos $n=0$ y $\alpha=1$.
\end{ejercicio}
